\scanpage{51}
nienne sur le corps des fractions \deleted{de} $L$ de $K$ ?

On notera que même dans le cas où $A=W(k)$, $k$ alg. clos, le
foncteur en question n'est pas \deleted{essentiellem} essentiellement surjectif.

Notons le

\uline{Problème} 5. \inserted{($A$ anneau de val. discrète complet d'inégales car., $k$ parf)} Expliciter le foncteur naturel qui va de la sous-catégorie pleine précédente de la catégorie des couples \deleted{d} $(M,\mathrm{Fil})$ d'un module
de Dieudonné $M$ sur \deleted{$W$} $k$ (i.e. \deleted{sur} module sur $W$) et d'une filtration à
deux crans de $M\otimes_W L$, vers les \deleted{mo}dules galoisiens sur $L$ (représent\unsure{a}tions
\inserted{vec\deleted{\unsure{toir}}torielles des dimension finie sur $\underline Q_p$ de}
de $\mathrm{Gal}(\bar L/L)$)\deleted{dans des \unsure{représen|espaces} modules libres de type fini} Peut-on même
définir ce foncteur \deleted{sur la sous-catégorie} de façon naturelle sans se
borner à cette sous-catégorie pleine,\deleted{\unsure{a}}en imposant seulement des conditions de rangs sur le gradué, mettons ?

\uline{Problème} 6. Revenant à un $A$ comme dans Pb 4, peut-on trouver un foncteur \deleted{d'une \unsure{c}} de la \deleted{\unsure{a}}catégorie des couples $(M,\mathrm{Fil}^{1})$ d'un isocristal
de Dieudonné \inserted{admissible} sur $A_{\unsure{o|0}}=A/pA$ (ou encore, sur $A_{o\ \mathrm{red}}$ \deleted{\illegible{1}}) et d'une filtration en deux crans par modules localement libres de $(M_A)_p$, ayant le
rang qu'il faut\deleted{,} (et sou\deleted{s}mis éventuellement à d'autres conditions
ouvertes \inserted{au sens Zar}, assurant que le couple envisagé provient d'un isogroupe de
BT sur $A$),vers les faisceaux $p$-adiques constants tordus sur $U=\mathrm{Spec}(A_p)$,
qui exprime le foncteur du théorè\deleted{\unsure{p}}me (VII 10.8.2) précédent.

\deleted{4) Question des relations entre cohomologie cristalline et $p$-adique.}

\deleted{\deleted{Problème} \uline{Conjecture} 7. On peut, pour tout anneau de valuation discrète complet $V$ d'inégales caractéristiques, de corps \deleted{des fractions} ris\deleted{\unsure{u}}\unsure{i|u}duel parfait (pour simplifier), de corps des fractions $L$, trouver un
foncteur \deleted{\unsure{san}}canoniques sur une sous-catégorie pleine convenable des couples $(M,\mathrm{Fil}^{\deleted{1}})$ \deleted{envisagée} dans le Pb 5 vers les}
% archive sheet number: 50

